\documentclass[a4paper,11pt]{ctexart} % 使用 ctexart 完美支持中文

% ==========================================
% 1. 宏包配置
% ==========================================
\usepackage[margin=1in]{geometry}
\usepackage{times}
\usepackage{graphicx}
\usepackage{amsmath}
\usepackage{xcolor}
\usepackage[ruled,vlined,linesnumbered]{algorithm2e}
\usepackage{caption}
\usepackage{pifont} 

\title{Debug: Syntax-Aware Query Abstraction}
\author{Morph Team}
\date{\today}

\begin{document}

\maketitle

\section{Approach}

\subsection{Syntax-Aware Query Abstraction}
\label{sec:abstraction}

% ============================================================
% 算法 2：规则引导的递归方言实例化
% ============================================================
\begin{algorithm}[htbp]
\small
\caption{规则引导的递归方言实例化}
\label{alg:translation}

% 仅将标签换成中文，保留控制流关键字为英文
\SetKwInOut{Input}{输入}
\SetKwInOut{Output}{输出}

\Input{AQT $\mathcal{T} = \langle S_{norm}, \mathcal{D}, \mathcal{G}, \mathcal{H}, \mathcal{C} \rangle$；目标方言 $D_t$；混合知识库 $HKB$；阈值 $\delta$}
\Output{兼容的目标查询种子 $Q_{target}$}

$Order \leftarrow \textsc{TopologicalSort}(\mathcal{G})$; $Fragments \leftarrow \emptyset$ \label{line:alg2:sort} \tcp*{自底向上的拓扑排序}
\ForEach{$id$ \textbf{in} $Order$}{
    $S \leftarrow \textsc{GetSegment}(\mathcal{T}, id)$\;
    \ForEach{$c$ \textbf{in} $\mathcal{G}.children(id)$}{
        $S \leftarrow \textsc{Substitute}(S, c, Fragments[c])$ \label{line:alg2:substitute_start} \tcp*{替换子查询占位符}
    } \label{line:alg2:substitute_end}
    \ForEach{$\tau$ \textbf{in} $\textsc{AbstractTokens}(S)$}{
        $(rule, conf) \leftarrow \textsc{Retrieve}(HKB, \tau, D_t)$ \label{line:alg2:retrieve} \tcp*{查询混合知识库}
        \If{$conf < \delta$}{
            $rule \leftarrow \textsc{Generator}(\tau, D_t)$ \label{line:alg2:generate} \tcp*{调用 LLM 生成转换规则}
            $HKB.\textsc{Cache}(\tau, rule)$\;
        }
        $frag \leftarrow \textsc{Translator}(\tau, rule, \mathcal{D})$ \label{line:alg2:translate} \tcp*{实例化 Token}
        \While{$err \leftarrow \textsc{Parse}(frag, D_t)$}{
            $frag \leftarrow \textsc{Fixer}(frag, err)$ \label{line:alg2:fix} \tcp*{错误修复循环}
        }
        $S \leftarrow \textsc{Replace}(S, \tau, frag)$\;
    }
    $Fragments[id] \leftarrow S$ \label{line:alg2:cache} \tcp*{缓存已实例化的片段}
}
\Return $Fragments[\textsc{Root}(\mathcal{G})]$ \label{line:alg2:assemble} \tcp*{组装并返回最终查询}
\end{algorithm}


% ============================================================
% 算法 2：规则引导的递归方言实例化
% ============================================================
\begin{algorithm}[htbp]
\small
\caption{规则引导的递归方言实例化}
\label{alg:translation}

% 仅将标签换成中文，保留控制流关键字为英文
\SetKwInOut{Input}{输入}
\SetKwInOut{Output}{输出}

\Input{AQT $\mathcal{T} = \langle S_{norm}, \mathcal{D}, \mathcal{G}, \mathcal{H}, \mathcal{C} \rangle$；目标方言 $D_t$；混合知识库 $HKB$；阈值 $\delta$}
\Output{兼容的目标查询种子 $Q_{target}$}

$Order \leftarrow \textsc{TopologicalSort}(\mathcal{G})$; $Fragments \leftarrow \emptyset$ \label{line:alg2:sort} \tcp*{自底向上的拓扑排序}
\ForEach{$id$ \textbf{in} $Order$}{
    $S \leftarrow \textsc{GetSegment}(\mathcal{T}, id)$\;
    \ForEach{$c$ \textbf{in} $\mathcal{G}.children(id)$}{
        $S \leftarrow \textsc{Substitute}(S, c, Fragments[c])$ \label{line:alg2:substitute_start} \tcp*{替换子查询占位符}
    } \label{line:alg2:substitute_end}
    \ForEach{$\tau$ \textbf{in} $\textsc{AbstractTokens}(S)$}{
        $(rule, conf) \leftarrow \textsc{Retrieve}(HKB, \tau, D_t)$ \label{line:alg2:retrieve} \tcp*{查询混合知识库}
        \If{$conf < \delta$}{
            $rule \leftarrow \textsc{Generator}(\tau, D_t)$ \label{line:alg2:generate} \tcp*{调用 LLM 生成转换规则}
            $HKB.\textsc{Cache}(\tau, rule)$\;
        }
        $frag \leftarrow \textsc{Translator}(\tau, rule, \mathcal{D})$ \label{line:alg2:translate} \tcp*{实例化 Token}
        \While{$err \leftarrow \textsc{Parse}(frag, D_t)$}{
            $frag \leftarrow \textsc{Fixer}(frag, err)$ \label{line:alg2:fix} \tcp*{错误修复循环}
        }
        $S \leftarrow \textsc{Replace}(S, \tau, frag)$\;
    }
    $Fragments[id] \leftarrow S$ \label{line:alg2:cache} \tcp*{缓存已实例化的片段}
}
\Return $Fragments[\textsc{Root}(\mathcal{G})]$ \label{line:alg2:assemble} \tcp*{组装并返回最终查询}
\end{algorithm}

\end{document}